\documentclass[11pt]{article}

\author{Math 1030}
\date{Due Thursday, April 23 by 5pm} 
\title{Homework 9}

\usepackage{graphicx,xypic}
\usepackage{amsthm}
\usepackage{amsmath,amssymb}
\usepackage{amsfonts}
\usepackage{xcolor}
\usepackage[margin=1in]{geometry}
\usepackage[shortlabels]{enumitem}
\newtheorem{problem}{Problem}
\renewcommand*{\proofname}{{\color{blue}Solution}}


\setlength{\parindent}{0pt}
\setlength{\parskip}{1.25ex}


\begin{document}

\maketitle

% DON'T LEAVE THIS BLANK! 
{\bf\Large Your Name:} 

% DON'T FORGET YOUR COLLABORATORS! 
Collaborator names: 

\vspace{.3in}
Topics covered: random graphs, spectral graph theory

Instructions {\bf (read these carefully!)}: 
\begin{itemize}
\item This assignment must be submitted on Gradescope by the due date. 
\item If you collaborate with other students (this is encouraged),  list your collaborators above. 
\item Homework is graded anonymously, so avoid putting your name on each page of the assignment.
\item If you are stuck, please ask for help (from me, Rafi, or a classmate). Use Ed discussions!  
\item You may freely use any fact proved in class (mention it in the appropriate spot). You \emph{may not} freely use facts from the book or other sources.
\item As a general rule, try to restrict your solution to each problem to a single page. Often solutions can be short, while still being complete. If your solution is very long, think more about how you might express it more concisely.
\end{itemize}
\pagebreak 


\begin{problem}
Prove the Rado graph $R$ is homogeneous: any isomorphism between subgraphs of $R$ can be extended to an isomorphism of $R$. 
\end{problem}

\begin{proof}

\end{proof}

\pagebreak 

\begin{problem}
Let $G=(V,E)$ be a finite graph. Count among all partitions of the vertex set $V=A\sqcup B$, the total number of $(A,B)$ edges. Show that the average number of $(A,B)$ edges (among all partitions) is $|E|/2$. Deduce that every graph has a bipartite subgraph with at least half of its edges. 
\end{problem}

\begin{proof}

\end{proof}


\pagebreak 

\begin{problem}
Fix $r\le n$ and $s\le m$. Count the total number of monochromatic copies of $K_{r,s}$ among all edge $2$-colorings of $K_{n,m}$. \footnote{First count the number of colorings where a given copy of $K_{r,s}\subset K_{n,m}$ is monochromatic. We did something similar in class.} Show that the average number of monochromatic $K_{r,s}$'s among all colorings is ${n\choose r}{m\choose s}2^{1-rs}$. 
\end{problem}

\begin{proof}

\end{proof}

\pagebreak 


\begin{problem}
Let $G$ be a finite graph, and write $A$ for the adjacency matrix. Using the approach from lecture, show that the largest eigenvalue of $A$ is equal to the maximum vertex degree if and only if $G$ is regular. 
\end{problem}

\begin{proof}

\end{proof}


\pagebreak 


\begin{problem}
Let $B$ be a matrix (not necessarily square) and consider 
\[A=\left(\begin{array}{cc}0&B\\B^t&0\end{array}\right).\]
Show that the eigenvalues of $A$ occur in $\pm$ pairs.\footnote{This is a purely linear algebra problem. Assuming $(x,y)$ is an eigenvector with eigenvalue $\lambda$ find an eigenvector with eigenvalue $-\lambda$.} 

Show that a graph $G$ is bipartite if and only if the eigenvalues of the adjacency matrix $A$ occur in $\pm$ pairs.\footnote{The first part of the problem gives one direction. For the other direction, use (1) the trace of a matrix is the sum of its eigenvalues, (2) eigenvalues of $A^k$ have the form $\lambda^k$ where $\lambda$ is an eigenvalue of $A$ (when $A$ is diagonalizable), and (3) the trace of $A^k$ count certain kinds of walks in $G$ (be specific).}

\end{problem}

\begin{proof}

\end{proof}




\end{document}